\documentclass[twocolumn]{article}
\usepackage{graphicx}
\usepackage{amsmath, amssymb, amsfonts}
\usepackage{geometry}
\usepackage{tikz}
\usetikzlibrary{arrows.meta, positioning, shapes.geometric, calc}
\usepackage{caption}
\usepackage{hyperref}
\usepackage{booktabs}
\usepackage{multirow}

\geometry{a4paper, left=20mm, right=20mm, top=25mm, bottom=25mm}

\title{\textbf{Discussion: Integrating Brain-Computer Interfaces with Adaptive AI}}
\author{Student ID: 11314389}
\date{\today}

\begin{document}

\maketitle

\section{Discussion}

\subsection{Addressing Signal Non-Stationarity via Spatiotemporal Sequence Modeling}

A primary barrier to robust motor imagery decoding is non-stationary signal drift, which degrades the performance of single-trial convolutional networks (CNNs) during continuous operation \cite{hameed2025enhancing, prev_student2024}. By formulating brain-computer interface (BCI) decoding as a sequential decision-making process, CTNet mitigates this drift. The integration of an LSTM module within the Q-network allows the model to retain past state representations, effectively arbitrating between immediate, potentially noisy sensory evidence and temporal context.

Empirically, this spatiotemporal modeling translates to a mean accuracy of 79.4\% (91.3\% peak) on the 4-class BCI Competition IV-2a dataset. In contrast, standard feed-forward architectures such as EEGNet and FBCSP inherently lack explicit temporal memory across epochs, typically plateauing near 68--75\% and 68--70\%, respectively. Our reproduced OVR-CSP+LDA baseline correspondingly achieved 74.3\%. The performance margin observed with CTNet supports the hypothesis that recurrent sequential modeling reduces susceptibility to short-term signal fluctuations, enabling more stable representations for downstream control.

\subsection{Evaluating Policy Optimization Against Feature Selection Frameworks}

Recent advancements in BCI have increasingly leveraged reinforcement learning (RL). Notably, multi-agent frameworks such as MARS \cite{shin2024mars} utilize RL to optimize spatial-spectral and temporal feature selection offline, reporting an 86.78\% subject-dependent accuracy on the IV-2a dataset. While CTNet's mean performance (79.4\%) is lower than this specialized feature-selection benchmark, its peak subject accuracy (91.3\%) demonstrates a highly competitive upper bound.

More importantly, the objective of CTNet differs fundamentally from feature-selection models like MARS. CTNet is designed for continuous closed-loop robotic actuation. Rather than optimizing the feature space offline, CTNet trains a target policy (via Deep Q-Network) directly upon interpretable spatial maps to output discrete end-effector commands \cite{nallani2024rleegnet}. This architecture directly maps neural patterns to action sequences, generating smooth, logically continuous trajectories that are critical for safe physical actuation---a capability that static classifiers cannot inherently provide.

\subsection{Generalization to High-Density Neural Recordings}

To validate the generalizability of the proposed framework, we evaluated CTNet on the large-scale PhysioNet EEG Motor Movement/Imagery Dataset (64 channels, 109 subjects). The results confirm that CTNet scales effectively to high-density EEG configurations without requiring architectural modifications.

Against modern 4-class motor imagery baselines---such as hybrid CNN-Transformer frameworks achieving 76.4\% on PhysioNet, and CNN1D\_MF models achieving 69.2\% on IV-2a---CTNet proves highly competitive. The successful deployment across both datasets demonstrates that the 1D-CNN+LSTM topology extracts robust, cross-dataset motor imagery features. This indicates that the performance gains are derived from the architecture's inherent spatiotemporal modeling rather than overfitting to the specific 22-channel layout of the IV-2a dataset.

\subsection{The Pragmatic Case for Binary Classification in Robotic Control}

While 4-class decoding serves as a rigorous benchmark, practical continuous control---such as planar target tracking via a robotic arm or wheelchair navigation---often benefits from reduced action spaces to minimize critical failure modes. When evaluated exclusively on lateral (left/right) commands, the OVR-CSP+LDA baseline accuracy improved significantly to 85\%.

This high-reliability binary classification validates the pipeline's readiness for physical deployment. In authentic environments, complex continuous trajectories are frequently synthesized from highly reliable, hierarchical binary decisions rather than error-prone 4-class direct mappings. Embedding these robust binary state transitions within the RL Markov Decision Process constitutes the core practical advantage of CTNet, ensuring the low error rates required for safe, reproducible robotic motion planning.

\begin{thebibliography}{9}
\bibitem{hameed2025enhancing}
I.~Hameed \emph{et al.}, ``Enhancing motor imagery EEG signal decoding through machine learning,'' \textit{Computers in Biology and Medicine}, 2025.

\bibitem{prev_student2024}
W.~Zhao \emph{et al.}, ``A Brain-Computer Interface Control System Design Based on Deep Learning,'' University of Manchester, 2024.

\bibitem{nallani2024rleegnet}
S.~Nallani and G.~Ramachandran, ``RLEEGNet: Integrating Brain-Computer Interfaces with Adaptive AI,'' \textit{arXiv:2402.09465}, 2024.

\bibitem{shin2024mars}
D.-H.~Shin \emph{et al.}, ``MARS: Multiagent reinforcement learning for spatial-spectral and temporal feature selection,'' \textit{IEEE Transactions on Systems, Man, and Cybernetics: Systems}, 2024.
\end{thebibliography}

\end{document}
